\chapter{Metric Definitions}\label{app:metrics}

The metrics were fixed before the first robot experiment and revised twice, each time with every result recomputed from
the raw files: version~2 (before the pilot run) counts a trial ended by a fault as a failure instead of repeating it;
version~3 (before setup~v2) reads the fault label from the run's own record, counts jerk only while the arm moves, and
adds the recovery time (M11) and the deviation from the policy's plan (M12). A \emph{tick} is one 30~Hz control step of
the robot computer; the \emph{first action} is the first tick that executes an action from a chunk; \emph{arm joints}
are the five joints other than the gripper, in degrees.

\begin{longtable}{@{}l >{\raggedright\arraybackslash}p{0.2\textwidth} p{0.66\textwidth}@{}}
\caption{Metric definitions (version 3).}\label{tab:metricsfull}\\
\toprule
& Metric & Definition \\
\midrule
\endfirsthead
\toprule
& Metric & Definition \\
\midrule
\endhead
\bottomrule
\endfoot
M1 & Success rate & The operator presses the success key before the time limit and the result meets the written
  success rule. A failure key, the time limit, or a terminal fault after the first action is a failure. Reported as
  successes over valid trials with a Wilson 95\% interval. \\
M1f & Ended by a fault & Trials that the control system ended after the first action (failures in M1), read from the
  run's outcome and first action whatever key the operator pressed. \\
M1t & Time to success & Seconds from the first action to the success key; successful trials only. \\
M2 & Stall fraction & Ticks after the first action at which no new action was available (the arm holds), over all
  ticks after the first action. \\
M3 & Stall events, longest stall & A stall event is a run of consecutive stall ticks; the longest stall is the longest
  run, in seconds. \\
M4 & Turnaround & For every accepted chunk, the time on the robot computer from sending the observation to receiving
  its chunk; per trial the median and the 95th percentile. \\
M5 & Latency estimate error & For every accepted chunk, the scheduler's latency estimate when the request was sent,
  minus the measured turnaround (positive: an overestimate). \\
M6 & Server inference time & For every accepted chunk, the server's reported time to produce it. M4 minus M6 is the
  network and queueing share. \\
M7 & Jerk & The commanded arm-joint positions on the tick grid (held through stalls), third finite difference times
  $30^3$, root mean square over the five arm joints and the samples while the arm moves (a sample spans four ticks
  and counts when at least one of them executed a new action), in \jerkunit. \\
M8a & Chunk switch disagreement & At every switch to a new chunk, over the steps both chunks cover, the mean Euclidean
  distance between their arm-joint actions. \\
M8b & Largest step & The largest change of any arm joint between two consecutive commanded positions. \\
M9 & Held-out action error & On the held-out episodes, from every 10th frame, the policy predicts a chunk from the
  recorded observation; the first 15 predicted actions are compared with the 15 recorded actions that follow. Mean
  absolute error: arm joints (mean of the five) and gripper. Flow and diffusion policies: one sample with a fixed
  seed. \\
M10 & Inference time, GPU memory & The server's inference time over 200 requests from held-out observations after 10
  warm-up requests (median and 95th percentile), and the peak GPU memory. \\
M11 & Recovery time & For each timed fault (spike: until its delayed reply is released; drops, outage, rise: until the
  window ends; with the trial's shift): seconds from the fault's end to the end of the first stall still going then or
  starting within 1~s after it, or 0 when there is none. Not recovered: the arm never moves again, or the system ended
  the run before the fault ended. Not counted for that fault: the operator's key or success ended the run first.
  Reported as recovered over counted runs and the median of the recovered. \\
M12 & Deviation from the plan & On every executed tick, the mean over the five arm joints of the absolute difference
  between the commanded position and the newest accepted chunk's raw action for that step (before any smoothing);
  0 without smoothing. \\
\end{longtable}

\paragraph{Reporting rules.}
\begin{itemize}
\item Ten valid trials per condition, policy and task (robot) or eight runs per cell (bench), in rounds with a seeded
      random order; plans may grow to more trials or more models, never change otherwise.
\item A trial voided for a setup mistake is listed but not counted; a fault before the first action is not counted
      and runs again.
\item Every metric is reduced to one value per trial before it is summarised (median and interquartile range).
\item Comparisons are made within one plan and for the same policy; each plan names its comparisons, and only its
      primary ones are tested, Holm-corrected over the plan; the others are exploratory.
\item A report never averages over different robot computers or servers; the timing metrics are shown per robot
      computer and server.
\item On the replay robot, M2--M8, M11 and M12 are reported, not M1; a bench run whose control loop ran below 27~Hz is
      not counted.
\end{itemize}
